\section{商业电站：资金、风险与路线选择}

前面讲装置，本章讲钱和路线。杨钊给出的路线分三台关键装置：洪荒 70 验证完整装置工程可行性；洪荒 170 以较低成本实现高参数和 Q 大于 10；洪荒 380 则作为完整示范电站，目标是形成可卖给电力业主的产品。170 大约是三十亿人民币级别投入，380 则是百亿到二百亿人民币量级的示范电站成本。

这里最重要的是风险转化。能量奇点选择相对保守的物理设计，尽量不冒新的科学风险；真正要解决的是高参数工程风险和成本风险。与 Helion 等其他路线相比，高温超导托卡马克的优势在于托卡马克路线已有大量高参数实验积累；而一些其他磁约束位形可能还需要把实验参数外推几个数量级，这在杨钊看来属于科学风险更高的路线。

\begin{warningbox}{科学风险与工程风险要分开}
科学风险是“答案是否存在、物理过程是否按外推工作”；工程风险是“答案存在，但能否做出满足参数、成本和可靠性的系统”。创业公司更适合压工程风险和成本风险；高科学风险问题更适合科研机构长期探索。
\end{warningbox}

\begin{knowledgebox}{资金与里程碑表}
\begin{tabularx}{\textwidth}{p{0.18\textwidth}p{0.32\textwidth}X}
\toprule
阶段 & 目标 & 资金/商业含义 \\
\midrule
洪荒 70 & 完整高温超导装置工程可行 & 证明新材料路线能建成、能运行。 \\
洪荒 170 & 低成本实现 Q 大于 10 & 需要大额融资，证明商业化核心装置可行。 \\
洪荒 380 & 示范电站，约 500MW 级产品设想 & 面向电力业主销售，验证电站经济性。 \\
规模复制 & 多台商业电站 & 度电成本继续下降，进入能源市场竞争。 \\
\bottomrule
\end{tabularx}
\end{knowledgebox}

\begin{figure}[H]
\centering
\includegraphics[width=0.92\textwidth]{holy-grail-tradeoff.png}
\caption{能源圣杯的代价：无限能源愿景背后是极端工程复杂度。自制概念图，依据 01:30:00--02:00:00 对谈内容整理。}
\end{figure}

\begin{knowledgebox}{读图：圣杯之所以难，是因为代价同时在五个方向出现}
高温、磁约束、材料寿命、资金消耗和时间不确定性共同构成聚变商业化难度。任何单点乐观都不够，必须把物理、工程、资本和监管放进同一张账本。
\end{knowledgebox}

\subsection{度电成本是金标准}

本节把前面所有技术路线收束到商业判断。Q 值、磁场、装置尺寸和子系统参数都很重要，但它们最终要服务一个朴素指标：发一度电要多少钱。只有当度电成本接近火电，聚变才开始进入能源系统；如果能低一个数量级，才可能改变能源使用方式。

访谈中杨钊反复强调，聚变商业化的金标准是度电成本。只要度电成本接近火电，聚变就有进入能源系统的可能；如果能比火电低一个数量级，能源供给格局和文明使用能源方式都会发生更深变化。第一台示范电站成本必然较高，但如果它已经接近火电，市场就能看到规模复制后的下降空间。

\begin{importantbox}{商业化不是第一度电，而是低成本第一度电}
发出第一度电是技术里程碑；按可接受造价、监管和运维成本稳定发电，才是商业里程碑。核聚变从示范到主流能源，中间还要跨过长时间运行、取热发电、维护、监管和供应链复制。
\end{importantbox}

\begin{knowledgebox}{三台装置的商业含义}
\begin{tabularx}{\textwidth}{p{0.18\textwidth}p{0.34\textwidth}X}
\toprule
装置 & 主要目标 & 商业解释 \\
\midrule
洪荒 70 & 验证完整高温超导托卡马克可建成、可运行 & 证明“钢船能下水”，不是直接证明电站经济性。 \\
洪荒 170 & 以较低成本实现 Q 大于 10 的实验装置 & 证明高温超导路线能把科学装置压成商业可行装置。 \\
洪荒 380 & 500MW 级示范电站设想，面向业主交付 & 检验长时间运行、取热发电和电站产品化。 \\
\bottomrule
\end{tabularx}
\end{knowledgebox}

\subsection{本章小结}

聚变创业的真正产品不是论文和装置照片，而是成本可接受的电站。洪荒路线的核心，是把科学风险尽量前置为已验证物理，把创业公司资源集中在工程交付、子系统验证和度电成本下降上。

